% GENERATED FILE. Edit canonical lessons, then run npm run generate:books.
% canonical-source-hash: fnv1a64:d1f12fcad7428f80
% canonical-lessons: SA-C532-upayanam, SA-C532-gucchah, SA-C532-dukulam, SA-C532-mekhala, SA-C532-upaharati

\chapter{Choosing a Gift}
\label{ch:sa-choosing-a-gift}
\hlchaptermodality{532}

\begin{chapteropening}
\textbf{By the end of this chapter:} \emph{I can choose a gift and say what I will buy if it is right.}

Complete the last lesson of chapter 532: I can choose a gift and say what I will buy if it is right.
\end{chapteropening}

\section[\texorpdfstring{\sk{उपायनम्} (upāyanam)}{उपायनम् (upāyanam)}]{\sk{उपायनम्} (upāyanam) — a present}

\label{lesson:SA-C532-upayanam}



\begin{quote}
Before the new one: say the Sanskrit for a prayer, then the Sanskrit for a desire.
\end{quote}

\subsection*{You'll want to know: \sk{उपायनम्}}
\textbf{\sk{उपायनम्}} — \emph{upāyanam} — \textquotedblleft{}a present\textquotedblright{}.

\begin{quote}
\sk{यदि} \sk{मित्रम्} \sk{आगच्छति}, \sk{तर्हि} \sk{अहं} \sk{तस्मै} \sk{उपायनं} \sk{ददामि।}
\end{quote}

\begin{quote}
\emph{yadi mitram āgacchati, tarhi ahaṃ tasmai upāyanaṃ dadāmi.}
\end{quote}

\begin{quote}
If my friend comes, I will give him a present.
\end{quote}

A neuter noun.

\begin{cousinweb}
\textbf{\sk{उपायनम्}} (\emph{upāyanam}) is built on \emph{upa}, \emph{ā} and the root √i, \textquotedblleft{}to go\textquotedblright{}: a coming near; MW gives \emph{upāyana} also as an offer, a present, a gift (MW).
\end{cousinweb}

\begin{grammarlens}[title={What you've built}]
One more word for saying what happens if.
\end{grammarlens}

\subsection*{Your turn}
\begin{itemize}
\raggedright
  \item \emph{Say it:} \emph{upāyanam}
  \item \emph{Say it:} \emph{upāyanam}, once more
  \item \emph{Say it:} \emph{yadi mitram āgacchati, tarhi ahaṃ tasmai upāyanaṃ dadāmi.} — If my friend comes, I will give him a present.
  \item \emph{Recall:} say the Sanskrit for a prayer, then the Sanskrit for a desire, then say \emph{upāyanam} again
\end{itemize}

\begin{tcolorbox}[breakable,colback=teal!4,colframe=teal!35!black,title={Before you move on}]
What does \sk{उपायनम्} mean? (\textquotedblleft{}A present\textquotedblright{}.) Say it once more.
\end{tcolorbox}
\section[\texorpdfstring{\sk{गुच्छः} (gucchaḥ)}{गुच्छः (gucchaḥ)}]{\sk{गुच्छः} (gucchaḥ) — a bunch of flowers}

\label{lesson:SA-C532-gucchah}



\begin{quote}
Before the new one: say the Sanskrit for a desire, then the Sanskrit for a present.
\end{quote}

\subsection*{You'll want to know: \sk{गुच्छः}}
\textbf{\sk{गुच्छः}} — \emph{gucchaḥ} — \textquotedblleft{}a bunch of flowers\textquotedblright{}.

\begin{quote}
\sk{यदि} \sk{पुष्पाणि} \sk{सुन्दराणि} \sk{सन्ति}, \sk{तर्हि} \sk{एकं} \sk{गुच्छं} \sk{क्रीणामि।}
\end{quote}

\begin{quote}
\emph{yadi puṣpāṇi sundarāṇi santi, tarhi ekaṃ gucchaṃ krīṇāmi.}
\end{quote}

\begin{quote}
If the flowers are pretty, I'll buy a bunch.
\end{quote}

A masculine noun.

\begin{cousinweb}
\textbf{\sk{गुच्छः}} (\emph{gucchaḥ}) is the word MW gives for a bush and for a bunch of flowers, a cluster of blossoms, the same as \emph{gutsa}; MW gives no further derivation, so its origin is uncertain (MW).
\end{cousinweb}

\begin{grammarlens}[title={What you've built}]
One more word for saying what happens if.
\end{grammarlens}

\subsection*{Your turn}
\begin{itemize}
\raggedright
  \item \emph{Say it:} \emph{gucchaḥ}
  \item \emph{Say it:} \emph{gucchaḥ}, once more
  \item \emph{Say it:} \emph{yadi puṣpāṇi sundarāṇi santi, tarhi ekaṃ gucchaṃ krīṇāmi.} — If the flowers are pretty, I'll buy a bunch.
  \item \emph{Recall:} say the Sanskrit for a desire, then the Sanskrit for a present, then say \emph{gucchaḥ} again
\end{itemize}

\begin{tcolorbox}[breakable,colback=teal!4,colframe=teal!35!black,title={Before you move on}]
What does \sk{गुच्छः} mean? (\textquotedblleft{}A bunch of flowers\textquotedblright{}.) Say it once more.
\end{tcolorbox}
\section[\texorpdfstring{\sk{दुकूलम्} (dukūlam)}{दुकूलम् (dukūlam)}]{\sk{दुकूलम्} (dukūlam) — fine cloth}

\label{lesson:SA-C532-dukulam}



\begin{quote}
Before the new one: say the Sanskrit for a present, then the Sanskrit for a bunch of flowers.
\end{quote}

\subsection*{You'll want to know: \sk{दुकूलम्}}
\textbf{\sk{दुकूलम्}} — \emph{dukūlam} — \textquotedblleft{}fine cloth\textquotedblright{}.

\begin{quote}
\sk{यदि} \sk{मात्रे} \sk{दुकूलं} \sk{रोचते}, \sk{तर्हि} \sk{अहं} \sk{तत्} \sk{क्रीणामि।}
\end{quote}

\begin{quote}
\emph{yadi mātre dukūlaṃ rocate, tarhi ahaṃ tat krīṇāmi.}
\end{quote}

\begin{quote}
If my mother likes the fine cloth, I'll buy it.
\end{quote}

A neuter noun.

\begin{cousinweb}
\textbf{\sk{दुकूलम्}} (\emph{dukūlam}) is the word MW gives for a plant and for the very fine cloth made from its inner bark; MW gives no derivation, so its origin is uncertain (MW).
\end{cousinweb}

\begin{grammarlens}[title={What you've built}]
One more word for saying what happens if.
\end{grammarlens}

\subsection*{Your turn}
\begin{itemize}
\raggedright
  \item \emph{Say it:} \emph{dukūlam}
  \item \emph{Say it:} \emph{dukūlam}, once more
  \item \emph{Say it:} \emph{yadi mātre dukūlaṃ rocate, tarhi ahaṃ tat krīṇāmi.} — If my mother likes the fine cloth, I'll buy it.
  \item \emph{Recall:} say the Sanskrit for a present, then the Sanskrit for a bunch of flowers, then say \emph{dukūlam} again
\end{itemize}

\begin{tcolorbox}[breakable,colback=teal!4,colframe=teal!35!black,title={Before you move on}]
What does \sk{दुकूलम्} mean? (\textquotedblleft{}Fine cloth\textquotedblright{}.) Say it once more.
\end{tcolorbox}
\section[\texorpdfstring{\sk{मेखला} (mekhalā)}{मेखला (mekhalā)}]{\sk{मेखला} (mekhalā) — a girdle}

\label{lesson:SA-C532-mekhala}



\begin{quote}
Before the new one: say the Sanskrit for a bunch of flowers, then the Sanskrit for fine cloth.
\end{quote}

\subsection*{You'll want to know: \sk{मेखला}}
\textbf{\sk{मेखला}} — \emph{mekhalā} — \textquotedblleft{}a girdle\textquotedblright{}.

\begin{quote}
\sk{यदि} \sk{एषा} \sk{मेखला} \sk{सुन्दरी} \sk{अस्ति}, \sk{तर्हि} \sk{ताम्} \sk{अहं} \sk{क्रीणामि।}
\end{quote}

\begin{quote}
\emph{yadi eṣā mekhalā sundarī asti, tarhi tām ahaṃ krīṇāmi.}
\end{quote}

\begin{quote}
If this girdle is pretty, I'll buy it.
\end{quote}

A feminine noun.

\begin{cousinweb}
\textbf{\sk{मेखला}} (\emph{mekhalā}) is the word MW records for a girdle, a belt; MW gives no derivation, so its origin is uncertain (MW).
\end{cousinweb}

\begin{grammarlens}[title={What you've built}]
One more word for saying what happens if.
\end{grammarlens}

\subsection*{Your turn}
\begin{itemize}
\raggedright
  \item \emph{Say it:} \emph{mekhalā}
  \item \emph{Say it:} \emph{mekhalā}, once more
  \item \emph{Say it:} \emph{yadi eṣā mekhalā sundarī asti, tarhi tām ahaṃ krīṇāmi.} — If this girdle is pretty, I'll buy it.
  \item \emph{Recall:} say the Sanskrit for a bunch of flowers, then the Sanskrit for fine cloth, then say \emph{mekhalā} again
\end{itemize}

\begin{tcolorbox}[breakable,colback=teal!4,colframe=teal!35!black,title={Before you move on}]
What does \sk{मेखला} mean? (\textquotedblleft{}A girdle\textquotedblright{}.) Say it once more.
\end{tcolorbox}
\section[\texorpdfstring{\sk{उपहरति} (upaharati)}{उपहरति (upaharati)}]{\sk{उपहरति} (upaharati) — presents; offers}

\label{lesson:SA-C532-upaharati}



\begin{quote}
Before the new one: say the Sanskrit for fine cloth, then the Sanskrit for a girdle.
\end{quote}

\subsection*{You'll want to know: \sk{उपहरति}}
\textbf{\sk{उपहरति}} — \emph{upaharati} — \textquotedblleft{}presents; offers\textquotedblright{}.

\begin{quote}
\sk{यदि} \sk{गुरुः} \sk{आगच्छति}, \sk{तर्हि} \sk{शिष्यः} \sk{तस्मै} \sk{फलानि} \sk{उपहरति।}
\end{quote}

\begin{quote}
\emph{yadi guruḥ āgacchati, tarhi śiṣyaḥ tasmai phalāni upaharati.}
\end{quote}

\begin{quote}
If the teacher comes, the pupil offers him fruit.
\end{quote}

Past (he …): \sk{सः} \sk{गुरवे} \sk{पुष्पाणि} \sk{उपहृतवान्}, \emph{saḥ gurave puṣpāṇi upahṛtavān}.

\begin{cousinweb}
\textbf{\sk{उपहरति}} (\emph{upaharati}) is built on \emph{upa}, \textquotedblleft{}near\textquotedblright{}, and the root √hṛ, \textquotedblleft{}to carry\textquotedblright{}, which also gives \emph{harati}; MW gives \emph{upahṛ} as to bring near, to offer, to present (MW).
\end{cousinweb}

\begin{grammarlens}[title={What you've built}]
One more word for saying what happens if.
\end{grammarlens}

\subsection*{Your turn}
\begin{itemize}
\raggedright
  \item \emph{Say it:} \emph{upaharati}
  \item \emph{Say it:} \emph{upaharati}, once more
  \item \emph{Say it:} \emph{yadi guruḥ āgacchati, tarhi śiṣyaḥ tasmai phalāni upaharati.} — If the teacher comes, the pupil offers him fruit.
  \item \emph{Recall:} say the Sanskrit for fine cloth, then the Sanskrit for a girdle, then say \emph{upaharati} again
\end{itemize}

\begin{tcolorbox}[breakable,colback=teal!4,colframe=teal!35!black,title={Before you move on}]
What does \sk{उपहरति} mean? (\textquotedblleft{}Presents; offers\textquotedblright{}.) Say it once more.
\end{tcolorbox}
